\scanpage{56}
\mspage{55}

c) $\varnothing\in\Phi$ \deleted{$\ne\varnothing$ [by virtue of b), this is always true if $ M\ne\varnothing$], and if $ M=\varnothing$, this means that $\varnothing$}

\deleted{\uline{\unsure{Rép.|Remark}} (is equivalent to : $\varnothing\in\Phi$)}

\uline{Remark} Conversely, let $ M$ be a set, $\deleted{\Phi}\mathcal F\subset\mathcal P( M)$ satisfying the conditions a) b) c). NB Condition c follows from \inserted{a),} b if $ M\ne\varnothing$. In all cases, it is equivalent to : \deleted{\unsure{$\in$|to}} c') $\deleted{\Phi}\mathcal F\ne\varnothing$ i.e there exists $\deleted{X}A\subset M$ such that $A\in\deleted{\Phi}\mathcal F$

% slanted
\margin{\unsure{These|The} conditions imply that $\mathcal F$ is stable under $A\cap B$ (but in general \unsure{we do not have|\ldots} $\not\Rightarrow$ $ M\in\mathcal F$ !)}

Then

1°) $\Phi$ ordered by inclusion \deleted{\unsure{ad|sa}}satisfies the conditions C1, C2. \deleted{The irreducible objects} \deleted{\unsure{are}}

2°) The irreducible objects are the $A_X$, and $X\mapsto\deleted{X}A_X$ is a bijection from $ M$ \deleted{with} the set of irreducible objects of $\deleted{\Phi}\mathcal F$.

3°) If $X\in M$, $A\in\deleted{\Phi}\mathcal F$, then $X\in A\Longleftrightarrow A_X\subset A$ \deleted{, hence} (tautological) hence \deleted{$A\in\mathcal F$} if we identify $ M$ \inserted{(by $X\unsure{\vdash|\mapsto} A_X$)} with the set of irreducible objects of $\deleted{\Phi}\mathcal F$, then for $A\in\mathcal F$, $\overline A=\{A_X\mid A_X\subset A\inserted{\text{ie.}\,X\in A}\}$ is identified with $A$ itself.

\centerline{\rule{1.8cm}{0.35pt}}

We shall consider $ M$ as ordered by $\leq$ (never by $\prec$ !) : multistrata and \uline{sub-multistrata}. The figures \inserted{$F$} can be identified with the \deleted{subsets} corresponding sets of multistrata $\overline F$ (\unsure{cl. by|each|each} figure with the set of its strata).
